\documentclass[12pt]{article}
\usepackage{amsfonts,amsthm,amsmath,amssymb,mathtools}
\usepackage{mathrsfs}
\usepackage{enumitem}
\usepackage{tikz-cd}
\usepackage{geometry}
\usepackage{mlmodern}

\geometry{letterpaper, portrait, margin=1in}

\setcounter{section}{0}

\renewcommand{\thesection}{\S\arabic{section}}
\renewcommand{\thesubsection}{\arabic{section}}
\newcommand{\x}{\mathbf{x}}
\DeclareMathOperator{\diam}{diam}
\newcommand{\dd}{\mathop{}\!\mathrm{d}}

\newtheorem{thm}{Theorem}
\newenvironment{theorem}[1]{
	\renewcommand{\thethm}{#1}
	\begin{thm}
		}{\end{thm}}

\newtheorem{lma}{Lemma}
\newenvironment{lemma}[1]{
	\renewcommand{\thelma}{#1}
	\begin{lma}
		}{\end{lma}}

\theoremstyle{definition}
\newtheorem{ex}{}
\newenvironment{exercise}[1]{
	\renewcommand{\theex}{#1}
	\begin{ex}
		}{\end{ex}}

\title{Homework 10}
\author{Connor Mooneyhan}
\date{April 9, 2026}

\begin{document}

\maketitle

\begin{ex}
	Let $c_0 = \left\lbrace (x_1, x_2, \dots) \in \ell^\infty : \lim_{n \to \infty} x_n = 0 \right\rbrace$.
	Show that $(c_0)'$ is isomorphic to $\ell^1$.
\end{ex}
\begin{proof}
	Given any $a = (a_n) \in \ell^1$, we have a linear functional $f_a$ which maps $x = (x_n) \in c_0$ by \begin{align*}
		f_a(x)  = \sum_{i = 1}^\infty a_ix_i ,
	\end{align*}
	which is bounded since $\left| \sum_{i=1}^\infty a_ix_i \right| \leq ||x||_{\infty}\sum_{i=1}^\infty a_i = ||x||_\infty||a||_1$.
	Now define $e_i$ to be the sequence with 1 in the $i^{\text{th}}$ position and 0 elsewhere.
	Then given a bounded linear functional $f \in (c_0)'$, we can define a sequence $a = (f(e_n))_{n \in \mathbb{Z}}$.
	Then given any $x = (x_1, x_2, \dots) \in c_0$, \begin{align*}
		f_a(x) & = \sum_{i=1}^\infty a_ix_i               \\
		       & = \sum_{i=1}^\infty f(e_i)x_i            \\
		       & = \sum_{i=1}^\infty f(x_ie_i)            \\
		       & = f\left(\sum_{i=1}^\infty x_ie_i\right) \\
           & = f(x),
	\end{align*}
  so $f$ is indeed realized by this mapping of elements of $\ell^1$ to functionals on $c_0$.
\end{proof}

\begin{ex}
	In an inner product space, we say that $x, y$ are orthogonal (sometimes we write $x \perp y$) if $\langle x, y \rangle = 0$.
	Show that if $x, y$ are non-zero and orthogonal then they must be linearly independent.
	Extend this result to an arbitrary finite number of pairwise orthogonal vectors and prove it.
\end{ex}
\begin{proof}
	I will just prove the general result since it includes the two-vector result.
	Suppose we have nonzero vectors $x_1, \dots, x_n$ which are all pairwise orthogonal.
	Then we want to show that \begin{equation}\label{1}
		c_1x_1 + \cdots + c_nx_n = 0
	\end{equation}
	implies $c_1 = \cdots = c_n = 0$.
	Indeed, assume we have \eqref{1} for some scalars $c_1, \dots, c_n$.
	Then for each $i \in \lbrace 1, \dots, n \rbrace$, \begin{align*}
		0 & = \left\langle c_1x_1 + \cdots + c_nx_n, x_i \right\rangle           \\
		  & = c_1\langle x_1, x_i \rangle + \cdots + c_n\langle x_n, x_i \rangle \\
		  & = c_i\langle x_i, x_i \rangle
	\end{align*}
	since $\langle x_j, x_i \rangle = 0$ whenever $i \neq j$ by hypothesis.
	Since $x_i$ is nonzero, $\langle x_i, x_i \rangle > 0$, so the final equality implies $c_i = 0$.
	Our choice of $i$ was arbitrary, so we are done.
\end{proof}

\begin{ex}
	Show that if $x_n$ converges to $x$ and $y_n$ converges to $y$ with respect to the induced norm then $\langle x_n, y_n \rangle$ converges to $\langle x, y \rangle$.
\end{ex}
\begin{proof}
	First, observe that \begin{align*}
		|\left\langle x_n, y_n \right\rangle - \left\langle x, y \right\rangle| & = |\left\langle x_n, y_n \right\rangle - \left\langle x, y_n \right\rangle + \left\langle x, y_n \right\rangle - \left\langle x, y \right\rangle| \\
		                                                                        & = \left\langle x_n - x, y_n \right\rangle + \left\langle x, y_n - y \right\rangle                                                                 \\
		                                                                        & \leq ||x_n-x||\cdot||y_n|| + ||x||\cdot||y_n-y||.
	\end{align*}
	We will study the final equation to get our result.
	The second term goes to zero because $||x||$ is constant with respect to $n$, and $||y_n - y||$ goes to zero as $n \to \infty$ since $y_n \to y$.
	The first term has $||x_n - x|| \to 0$ for the same reason, and of course $y_n \to y$, so that first term will also tend to 0 as $n \to \infty$.
	Thus we get that the entire RHS of that inequality goes to 0 as $n \to \infty$, and thus $|\langle x_n, y_n \rangle - \langle x, y \rangle|$ does as well since it is bounded above by 0.
\end{proof}

\begin{ex}
	Let $T : X \to X$ be a bounded linear operator on a complex inner product space $X$.
	If $\langle T(x), x \rangle = 0$ for all $x \in X$, show that $T = 0$.
	Also, give an example where this is false in real inner product spaces.
\end{ex}
\begin{proof}
	Let $a \in C$ and $u, v \in X$.
	Then \begin{align*}
		0 & = \left\langle T(u + av), u + av \right\rangle                                                                                                                       \\
		  & = \left\langle T(u) + aT(v), u + av \right\rangle                                                                                                                    \\
		  & = \left\langle T(u), u \right\rangle + \bar{a}\left\langle T(u), v \right\rangle + a \left\langle T(v), u \right\rangle + a\bar{a}\left\langle T(v), v \right\rangle \\
		  & = \bar{a}\left\langle T(u), v \right\rangle + a \left\langle T(v), u \right\rangle.
	\end{align*}
	Now since $a$ is arbitrary, we can look at a couple values of $a$ in particular.
	Let $A = \langle T(u), v \rangle$ and $B = \langle T(v), u \rangle$ for notational ease.
	Then when $a = 1$ we have \begin{equation*}
		A + B = 0
	\end{equation*}
	and when $a = i$ we have \begin{equation*}
		-iA + iB = 0 \implies -A + B = 0.
	\end{equation*}
	Adding the two equations together, we get that $2B = 0$, so $B = 0$, and then plugging that into the first equation we get $A = 0$.
	Therefore $\langle T(u), v \rangle = 0$ for all $u, v \in X$.
	If we choose $v = T(u)$, this yields $\langle T(u), T(u) \rangle = 0$, which is true only when $T(u) = 0$.
	This is true for all $u$, so $T = 0$.

	For the real counterexample, choose $T : \mathbb{R}^2 \to \mathbb{R}^2$ defined by $T(x, y) = (-y, x)$.
	Then indeed for any $(x, y) \in \mathbb{R}^2$, we have \begin{align*}
		\langle T(x, y), (x, y) \rangle & = \langle (-y, x), (x, y) \rangle \\
		                                & = -xy + xy                        \\
		                                & = 0.
	\end{align*}
	Thus $T$ satisfies the criterion, but it is clearly not the zero map (e.g. $(1, 0) \mapsto (0, 1)$).
\end{proof}

\end{document}
